\chapter{Arthrocentesis}

\section*{Overview}
Arthrocentesis is the aspiration of synovial fluid from a joint using a needle, performed for diagnosis, symptom relief, and intra-articular therapy.

\section*{Alternative Names}
Joint aspiration; joint tap; knee tap; intra-articular injection

\section*{Principle}
A needle is passed through the skin and joint capsule into the synovial space. Fluid is withdrawn for analysis (cell count, crystals, culture), and medication such as corticosteroid or local anesthetic can be instilled.

\section*{History}
Joint aspiration has been practiced since the 19th century for diagnosis of joint effusions. Aspiration became a cornerstone of gout, septic arthritis, and inflammatory arthritis diagnosis.

\section*{Why This Test or Procedure Is Ordered}
Performed to evaluate an acute monoarticular effusion (to exclude septic arthritis), diagnose crystal arthropathies, relieve large symptomatic effusions, inject corticosteroids or viscosupplements, and monitor therapy in inflammatory arthritis.

\section*{Is This a Primary or Secondary Test or Procedure}
Primary diagnostic and therapeutic procedure for joint effusions; a screening and monitoring tool in rheumatic disease.

\section*{How This Test or Procedure Is Performed}
The joint is identified and prepped sterilely, the skin and capsule may be anesthetized with lidocaine, and a needle (often 18-20 gauge) is advanced into the joint space. Fluid is aspirated, a portion is sent for analysis, and medication may be injected. The knee is the most common site.

\section*{What Preparation Is Needed}
Usually none beyond a targeted history and examination. Aspirin and anticoagulants are generally continued, but the physician assesses bleeding risk; septic joint precautions apply.

\section*{Contraindications and Precautions}
Absolute contraindications: overlying skin infection or cellulitis. Relative precautions: severe coagulopathy or anticoagulation, and a prosthetic joint (higher concern for introduced infection).

\section*{Risks}
Pain, bleeding, infection (rare with sterile technique), damage to cartilage or nerves, and post-injection flare.

\section*{What Happens After the Test or Procedure}
Pressure dressing, ice, and rest for 24-48 hours. Analgesia as needed. Results guide further workup.

\section*{Understanding Your Results}
Cloudy fluid suggests infection or inflammation; high white count with neutrophils suggests septic arthritis; monosodium urate crystals confirm gout; calcium pyrophosphate crystals confirm pseudogout; and a positive culture confirms septic arthritis.

\section*{Normal Results}
Clear or slightly yellow, low-viscosity fluid with low white blood cell count (usually under 200/$\mu$L) and no crystals or organisms.

\section*{Follow-up Tests and Procedures}
Culture and sensitivity if infection is suspected; repeat aspiration or imaging may be needed. Response to intra-articular injection is assessed clinically.

\section*{References}
\begin{enumerate}
  \item MedlinePlus. Joint aspiration. U.S. National Library of Medicine; 2025.
  \item Current Medical Diagnosis and Treatment. McGraw-Hill; 2025.
\end{enumerate}

\clearpage
